\section{问题一的模型建立与求解}\label{sec:model1}

\subsection{模型建立}

设方式1的位置观测为 $\{(t_k^{(1)}, x_k^{(1)}, y_k^{(1)})\}_{k=1}^{N_1}$（4\,Hz），方式2为 $\{(t_j^{(2)}, x_j^{(2)}, y_j^{(2)})\}_{j=1}^{N_2}$（5\,Hz）。附件1无噪声，两路观测同一真实轨迹 $\boldsymbol{p}(t)$，仅存在设备开机时间偏差 $\Delta\tau$：
\begin{equation}
    t_{\text{real}}^{(2)} = t^{(2)}_{\text{stamp}} + \Delta\tau
    \label{eq:time_offset}
\end{equation}

对两路分别建三次样条插值 $S_1(t)$、$S_2(t)$，则在对齐后公共时段 $\mathcal{T} = [t_{\max}, t_{\min}]$ 上，最小化残差平方和：
\begin{equation}
    \Delta\tau^* = \arg\min_{\Delta\tau} \sum_{t_k \in \mathcal{T}} \left\| S_1(t_k) - S_2(t_k - \Delta\tau) \right\|^2
    \label{eq:p1_obj}
\end{equation}

\subsection{算法设计}

采用\textbf{粗搜+精搜}两阶段策略：
\begin{enumerate}
    \item \textbf{粗搜}：在 $\Delta\tau \in [\Delta\tau_0 - 50, \Delta\tau_0 + 50]$（步长 0.1\,s）上枚举，取 $J$ 最小者为初始值。其中 $\Delta\tau_0 = t_1^{(1)}_{\text{start}} - t_1^{(2)}_{\text{start}}$。
    \item \textbf{精搜}：在粗搜最优 $\pm 1$\,s 范围内用 \texttt{scipy.optimize.minimize\_scalar}（bounded 方法，精度 $10^{-6}$\,s）细化。
\end{enumerate}

对齐完成后，10\,Hz 联合轨迹取两路样条均值（无噪声时两路重合）：
\begin{equation}
    \hat{\boldsymbol{p}}(t_k) = \frac{S_1(t_k) + S_2(t_k - \Delta\tau^*)}{2}, \quad t_k = t_0 + 0.1k
\end{equation}

\begin{algorithm}[H]
    \caption{问题一：时间对齐算法}
    \label{alg:solve1}
    \KwIn{两路观测 $\mathcal{D}_1, \mathcal{D}_2$}
    \KwOut{$\Delta\tau^*$, 10\,Hz 轨迹}
    分别建三次样条 $S_1, S_2$\;
    $\Delta\tau_0 \leftarrow t^{(1)}_{\text{start}} - t^{(2)}_{\text{start}}$\;
    \For{$\Delta\tau \in [\Delta\tau_0 - 50, \Delta\tau_0 + 50]$, step 0.1}{
        计算 $J(\Delta\tau)$，记录最小值\;
    }
    在最优 $\pm 1$\,s 内精搜 $\Delta\tau^*$\;
    在 10\,Hz 网格上输出均值轨迹\;
    \Return{$\Delta\tau^*$, 轨迹}\;
\end{algorithm}

\subsection{结果与分析}

\begin{table}[H]
    \centering
    \caption{问题一求解结果}
    \label{tab:result1}
    \begin{tabular}{@{}lrl@{}}
        \toprule
        指标 & 数值 & 说明 \\
        \midrule
        $\Delta\tau^*$ & $-198.4317$\,s & 方式2时间戳 + $\Delta\tau$ = 真实时间 \\
        对齐 RMSE & $8.39 \times 10^{-11}$\,m & 无噪声理论应为 $\sim 0$ \\
        公共时段长度 & 700.35\,s & $[270.4, 970.75]$ \\
        10\,Hz 轨迹点数 & 7004 & \\
        \bottomrule
    \end{tabular}
\end{table}

RMSE 达到 $10^{-11}$ 量级，验证了附件1确实无噪声，且三次样条 + 一维搜索方法能精确恢复时间偏差。算法复杂度为 $O(N \log N)$（样条构建）$+ O(M)$（网格搜索），远优于 DTW 的 $O(N^2)$。
